\chapter{Cosa cacciavano, epoca per epoca}
\label{cap:cosa_cacciavano}

Il capitolo precedente ha descritto, in generale, come si comportavano le specie animali che popolavano le Alpi e come la loro composizione sia cambiata nel tempo. Questo capitolo scende nel dettaglio, portando i dati concreti raccolti negli scavi dei siti prealpini veneti: cosa mangiavano davvero i cacciatori che li hanno frequentati, fase dopo fase, con le percentuali pubblicate dagli archeozoologi che hanno studiato quei resti, non con formule generiche come ``prevalentemente'' o ``soprattutto''.

\section{Le fasi più antiche: il Paleolitico medio in Lessinia}

Nei siti più antichi delle Prealpi venete, riferibili a un periodo compreso tra circa cinquantottomila e trentottomila anni prima dell'anno zero, attribuiti ai gruppi neanderthaliani, il quadro faunistico è dominato dagli ungulati di taglia media e grande.

\subsection{Riparo Tagliente}

Il Riparo Tagliente, ai piedi del Monte Tregnago, sul fianco sinistro della Valpantena, in Lessinia, è tra i più importanti siti paleolitici d'Italia: scoperto nel 1958 da Francesco Tagliente, appassionato di preistoria locale, e scavato dal 1962, prima dal Museo Civico di Storia Naturale di Verona e poi, dal 1967, dall'Università di Ferrara. La prima frequentazione, attribuita a Neanderthal, risale a un periodo compreso tra cinquantottomila e trentottomila anni prima dell'anno zero, verso la fine del Paleolitico medio.

Nei depositi musteriani più antichi (tagli da 52 a 47) l'associazione dei micromammiferi indica un ambiente aperto di prateria arborata; nei tagli successivi (dal 47 al 40) compaiono specie tipiche delle steppe asiatiche, segno di un clima continentale con inverni aridi e molto freddi; al tetto della sequenza musteriana (tagli 39-31) i resti indicano invece un clima più umido e temperato. La posizione geografica del sito, all'incrocio tra ambiente di pianura e ambiente montano, permise ai Neanderthal l'accesso a una notevole varietà di specie: sono diffuse le specie di ambiente forestale (soprattutto capriolo e cervo, più raramente cinghiale) e, in misura altrettanto consistente, le specie tipiche di montagna, in particolare camoscio, stambecco e marmotta \citep{capuzzisala1980}.

Un dato quantitativo utile per il confronto tra siti riguarda il cinghiale, la cui presenza percentuale varia notevolmente da un sito all'altro dell'Italia nord-orientale: allo stesso Riparo Tagliente raggiunge l'11,7\%, la percentuale più alta tra i siti tardoglaciali censiti nella regione \citep{fioreetal2003}.

\subsection{Altri siti lessinei del Paleolitico medio e superiore}

Un quadro più ampio, che comprende oltre a Riparo Tagliente anche Riparo Mezzena, Grotta A di Veja e Grotta della Ghiacciaia, tutti in Lessinia, mostra un ordine di abbondanza costante tra i siti: i cervidi di taglia media e piccola (cervo e capriolo) sono i taxa di caccia preferiti, probabilmente per la loro maggiore disponibilità nei pressi dei siti; seguono, in ordine di abbondanza, i bovidi di grande taglia (l'uro, \textit{Bos primigenius}, e il bisonte delle steppe, \textit{Bison priscus}) e quelli di taglia media (stambecco e camoscio) \citep{peresani2017ghiacciaia}. In misura molto minore e sporadica compaiono anche l'alce e il cervo gigante (\textit{Megaloceros giganteus}), meno rappresentati dei cervidi di taglia inferiore, oltre al cinghiale e, solo a Grotta di Fumane e a Riparo Tagliente, al cavallo selvatico.

\section{Le prime comunità di quota: l'Epigravettiano}

Con il miglioramento climatico descritto nel capitolo precedente, a partire da circa dodicimilasettecento anni prima dell'anno zero, il quadro comincia a cambiare. Gli insediamenti si spostano progressivamente verso quote più elevate, e la fauna cacciata si restringe a un numero più ridotto di specie, dominate nettamente da stambecco e camoscio.

\subsection{Riparo Dalmeri: la specializzazione più estrema}

Il Riparo Dalmeri, sull'Altopiano dei Sette Comuni, a milleduecentoquaranta metri di quota, rappresenta il caso di specializzazione venatoria più netto tra tutti i siti considerati in questo libro: lo stambecco costituisce da solo più del 90\% dei resti di mammiferi identificati nel livello 26c, il meglio documentato del sito. Le catture di capriolo, camoscio, alce e cinghiale, per contro, risultano del tutto occasionali. Tra i carnivori, l'orso è il più rappresentato, seguito da volpe e lupo, tutti documentati anche da resti di individui giovanili; il tasso è rappresentato da un solo individuo adulto. La presenza di resti di carnivori in età giovanile o neonatale, insieme a ossa che mostrano tracce di morsi carnivori, indica che il riparo fu usato anche come tana da questi animali, oltre che come sito di caccia umana; esistono però anche tracce sicure di macellazione umana su resti di orso e tasso, segno che i cacciatori sfruttavano anche questi carnivori quando l'occasione si presentava.

Un dato interessante, che distingue Dalmeri dagli altri siti epigravettiani, riguarda l'età degli stambecchi catturati: mentre nella maggior parte dei siti la caccia era diretta prevalentemente verso individui adulti (probabilmente per massimizzare la resa in carne), a Dalmeri gli stambecchi furono catturati in prevalenza in età giovanile \citep{fioretagliacozzo2005}.

\subsection{Riparo Soman: la caccia ai branchi di femmine}

Il Riparo Soman, anch'esso in ambiente prealpino veneto, offre un quadro diverso. Qui il camoscio è la specie prevalente, con percentuali superiori al 40\% dei resti, una delle percentuali più alte registrate in Italia per questa specie: su oltre trenta siti epigravettiani italiani analizzati in uno studio dedicato, solo Riparo Soman e tre siti dell'Italia centrale (Grotta Maritza, Grotta di Ortucchio, Grotta di Pozzo) superano questa soglia \citep{deangelis2021}. A Soman venivano cacciati anche stambecco, cervo, cinghiale, capriolo e alce, ma il camoscio resta la preda dominante, catturata in un periodo di frequentazione compreso tra l'estate e l'autunno.

Gli individui abbattuti a Soman comprendono animali di tutte le età, femmine adulte e piccoli inclusi: un pattern che suggerisce che la caccia si rivolgesse ai branchi di femmine con i loro piccoli, che d'estate si concentrano nelle praterie di alta quota. Nella seconda fase di occupazione del sito, i cervi catturati risultano in ugual misura adulti e giovani-adulti, un dettaglio che differenzia ulteriormente la strategia venatoria di questo sito da quella osservata altrove.

\subsection{Un confronto tra cinque siti}

Uno studio che ha confrontato sistematicamente cinque siti epigravettiani dell'Italia nordorientale, a quote e in ambienti diversi, offre una conclusione importante per il metodo di questo libro: le differenze osservate nei rapporti tra le specie di ungulati non sembrano dovute soltanto all'attività di caccia umana, ma riflettono e sono influenzate soprattutto dai cambiamenti climatici e dall'ambiente circostante ciascun sito \citep{fioreetal2003}. In altri termini: i cacciatori non sceglievano liberamente cosa cacciare, ma si adattavano a cosa il territorio circostante offriva loro, stagione dopo stagione. Lo stesso studio conferma il pattern già descritto: la caccia era diretta in prevalenza verso individui adulti, con le sole eccezioni degli stambecchi giovanili di Dalmeri e dei cervi di Soman, cacciati in ugual misura adulti e giovani-adulti.

\section{Uno sguardo d'insieme}

La tabella \ref{tab:fauna_siti} raccoglie i dati quantitativi discussi in questo capitolo, per permettere un confronto diretto tra i siti.

\begin{table}[htbp]
\centering
\small
\begin{tabular}{@{}p{2.6cm}p{2.0cm}p{2.6cm}p{1.4cm}p{2.4cm}@{}}
\toprule
\textbf{Sito} & \textbf{Quota} & \textbf{Specie dominante} & \textbf{Cinghiale} & \textbf{Note} \\
\midrule
Riparo Tagliente & fondovalle & mix cervidi/caprini di quota & 11,7\% & Paleolitico medio, Neanderthal \\
Riparo Dalmeri & 1.240 m & stambecco & 0,2\% & $>$90\% dei resti; catture giovanili \\
Riparo Soman & non specificata in questa sede & camoscio & 3\% & $>$40\%; caccia a branchi di femmine e piccoli \\
Riparo Villabruna A & non specificata in questa sede & non dominante singola & 4,1\% & confrontato negli studi di sintesi \\
\bottomrule
\end{tabular}
\caption{Confronto tra i principali siti archeozoologici prealpini veneti discussi in questo capitolo. I dati sul cinghiale, comparabili tra siti, sono tratti da Fiore, Tagliacozzo e Cassoli (2003).}
\label{tab:fauna_siti}
\end{table}

\section{Perché queste differenze non sono casuali}

Il confronto tra i siti discussi in questo capitolo conferma quanto anticipato nel capitolo precedente: non ha senso cercare, con lo stesso criterio, tracce di frequentazione preistorica in una valle di fondovalle ampio e in un valico di alta quota. Il Riparo Tagliente, in fondovalle, mostra un quadro faunistico misto, con una presenza significativa di cinghiale e di specie di ambiente aperto; Riparo Dalmeri, a milleduecentoquaranta metri, mostra una specializzazione quasi esclusiva sullo stambecco; Riparo Soman mostra una specializzazione altrettanto marcata, ma sul camoscio, con una strategia di caccia diversa (branchi di femmine e piccoli, non individui isolati).

Un criterio di ricerca che ignori questa distinzione rischia di cercare lo stesso tipo di segnale ovunque, quando in realtà il segnale atteso cambia con l'ambiente, con la quota e con l'epoca. È esattamente il tipo di dettaglio che una sintesi generica tende a perdere, e che invece risulta prezioso per chi debba costruire, come si farà nella terza parte di questo libro, un criterio di ricerca fondato su basi solide invece che su impressioni generiche.

Prima però, per completare il quadro necessario a quel lavoro, resta da affrontare un'ultima domanda, di natura diversa: cosa possono insegnare, su questi temi, i territori dove l'arte rupestre è invece abbondantissima, a partire proprio dalla Región de Atacama da cui questo libro è partito.
